\section{Discussion}

At model level, the comparison appears straightforward. On the same seven held-out records, the Morlet CNN--Conformer led the three study comparators in accuracy, macro-F1, and weighted-F1. Macro-F1 remained 0.26, however, and the class report changes how that lead should be read. F1 reached 0.74 for Normal and 0.57 for VEB, but 0.00 for SVEB and 0.01 for Fusion. The aggregate gain was carried by the two common categories with materially higher class scores; dependable recognition did not extend across all five classes.

Why does weighted-F1 reach 0.68 while macro-F1 remains 0.26? Normal and VEB account for nearly all held-out beats, so their scores dominate the weighted average. Macro-F1 gives the three sparse classes equal standing and exposes the gap in coverage. Reporting both measures prevents strong performance on the majority rhythms from being mistaken for broad class competence.

The seven test records were absent from model development. The resulting estimate therefore reflects a change in patient morphology and recording context rather than another sample of beats from familiar records. This stronger evaluation boundary comes with visible uncertainty: across the 38 development records, five-fold accuracy was $0.44\pm0.16$. We read that dispersion as evidence that patient composition materially affects the measured performance. A single small held-out cohort cannot show how much of the reported gain would survive another mix of records.

Architecture provides only a partial explanation. The attention-only and CNN--LSTM alternatives differed from the proposed model in more than one isolated module, so the comparison cannot attribute the gain to a single operation. Class-specific AUCs also resist a simple claim of dominance: CNN--LSTM had a slightly higher VEB AUC, and several alternatives exceeded the proposed model for Fusion or Q-class AUC. Within the implemented pipeline, local convolution followed by Conformer sequence modelling gave the best balance of accuracy, macro-F1, and weighted-F1. Explaining that result requires a component-wise ablation.

Of the 214 SVEB beats, 71\% were assigned to Normal and another 24\% to VEB. Among Fusion beats, 90\% were assigned to VEB. These directions are more informative for model development than the near-zero F1-scores alone. The SVEB-to-Normal pattern is compatible with a representation that relies more heavily on local morphology than on the timing cues available in a three-beat sequence. Fusion-to-VEB errors, in turn, suggest that ventricular-like local structure dominates the learned representation. The present analysis cannot distinguish these mechanisms. Holding the patient split fixed while varying sequence length, lead information, and the convolutional and attention components would test them directly.

Could the near-zero F1-scores be mainly a threshold problem? The post-hoc sweeps argue against that explanation. The best SVEB setting raised recall to 0.90 but reduced precision to 0.03; for Fusion, recall 1.00 was paired with precision 0.02. Such operating points would produce many false alarms for every correctly identified minority beat. Because the sweeps used the held-out probabilities, they cannot define a deployment threshold. Their shape nevertheless shows extensive overlap between rare-class outputs and competing classes, leaving little room for threshold selection alone to restore a useful precision--recall balance.

Comparison with published ECG classifiers cannot resolve the remaining uncertainty. Convolutional, recurrent, residual, and patient-adapted models operate on different signal durations and validation designs\cite{hannun2019cardiologist,kiranyaz2015personalized,luz2016ecg,acharya2017deep}. Even within the \textit{Journal of Electrocardiology}, CLINet performs signal-level classification whereas recent work on chronic heart failure uses long-term RR dynamics\cite{mantravadi2024clinet,pukkila2026chronic}. Each endpoint places different demands on representation, which makes direct ranking by headline scores uninformative. Here the useful comparison is internal: four implementations share one test cohort, and their errors can be inspected against the same class support.

Because the raw probabilities are archived, the confidence analysis can be reproduced alongside the class errors. Before scaling, Normal ECE was 0.417 and VEB ECE was 0.089; fitting one global temperature changed negative log-likelihood from 1.047989 to 1.047731. Recent work treats certainty and calibrated uncertainty as distinct outputs of AI-enabled ECG systems\cite{demolder2024certainty,doggart2025singlelead}. In this study, however, the temperature was fitted to the test cohort itself. The calculation is an internal diagnostic, and the negligible change in negative log-likelihood gives little reason to expect a global confidence adjustment to correct the class-specific recognition failure.

\subsection*{Limitations}

The evidence comes from one legacy public database and seven held-out records. This is a useful controlled test, but a small record count cannot represent population, device, or acquisition variability. Class support narrows the inference further: the test cohort contained 214 SVEB beats, 383 Fusion beats, and only two Q beats, making the Q results descriptive and the other rare-class estimates sensitive to a small number of patients. Beat-level bootstrap intervals quantify uncertainty among the observed beats; they do not measure uncertainty across new cohorts.

The input contains only the first ECG channel and three adjacent beats. Longer rhythm context or another lead could change the specific error pathways discussed above, and the current experiment cannot determine which omission matters more. Class-weighted cross-entropy was the only definitive imbalance treatment; no matched comparison was made with focal, cost-sensitive, or representation-learning alternatives. Finally, both threshold selection and temperature fitting used the held-out predictions after the definitive evaluation. They describe this model's behaviour on this cohort and should not be treated as independently validated operating procedures.

\subsection*{Priorities for further work}

An error-linked representation experiment should come first. With the record split and evaluation code held constant, longer temporal windows, multi-lead input, and targeted ablations can be judged by their effect on the SVEB-to-Normal and Fusion-to-VEB pathways. This would show whether the failures arise mainly from missing context, the learned local representation, or their interaction. Minority-class learning is the next question. Focal or explicitly cost-sensitive objectives, physiologically constrained augmentation, and self-supervised pretraining should be compared on development data, with operating points fixed before the untouched test cohort is examined. A successful method must improve the precision--recall balance for SVEB and Fusion; another increase in weighted-F1 would not answer the problem found here.

Only after those two questions are resolved should a fixed model move to external validation. Contemporary clinical and wearable cohorts would test whether the observed Normal/VEB discrimination survives changes in population, acquisition, and clinical endpoint, extending beyond the distinct sinus-rhythm prediction and 12-lead diagnostic tasks already reported\cite{attia2019artificial,ribeiro2020automatic}. Validation should include the preprocessing routes likely to appear in practice, including digitised ECGs and automated 12-lead interpretation systems, for which data-quality dependencies have been documented\cite{herman2024validation,demolder2025digitization}. If rare-class performance and independently selected confidence estimates remain acceptable there, a prospective clinician-in-the-loop study could then examine review efficiency, prioritisation, and false-alarm burden.
